%======================================================================
% ANEXOS TECNICOS - SANDRA DESKTOP CONTAINER
%======================================================================

\chapter{Anexos Tecnicos}
\label{ch:anexos}

Este apendice contiene informacion tecnica detallada sobre la arquitectura, flujos de datos y referencias de \textbf{Sandra Desktop Container}.

\section{Arquitectura de Software SDC}
\label{sec:arquitectura}

La arquitectura de SDC esta organizada en capas:

\begin{table}[H]
\centering
\begin{tabular}{p{4cm}p{5cm}p{5cm}}
\toprule
\textbf{Capa} & \textbf{Componentes} & \textbf{Tecnologia} \\
\midrule
Presentacion & Angular UI, Signals, SDC Services & Angular 15+ \\
\midrule
Puente & IPC Bridge, Commands API, Events API & Tauri 2.0 \\
\midrule
Core & Crypto, SQLite, WebSocket, Sandbox & Rust + Tokio \\
\midrule
Sistema & File System, Network, Process Mgmt & OS Nativo \\
\bottomrule
\end{tabular}
\caption{Arquitectura multicapa de SDC}
\label{tab:arquitectura-sdc}
\end{table}

\section{Ciclo de Vida de un Contenedor}
\label{sec:ciclo-vida}

Los contenedores en SDC siguen una maquina de estados:

\begin{enumerate}
    \item \textbf{CREADO}: El contenedor ha sido configurado pero no iniciado
    \item \textbf{CONFIGURADO}: Parametros de ejecucion establecidos
    \item \textbf{EJECUTANDO}: El contenedor esta activo y procesando
    \item \textbf{PAUSADO}: Ejecucion temporalmente detenida
    \item \textbf{DETENIDO}: Ejecucion finalizada, recursos liberados
    \item \textbf{DESTRUIDO}: Contenedor eliminado permanentemente
\end{enumerate}

\section{Flujo de Cifrado de Datos}
\label{sec:flujo-cifrado}

El proceso de cifrado en SDC:

\begin{enumerate}
    \item \textbf{Datos en Texto Plano} ingresan al sistema
    \item \textbf{Generacion de Clave}: Se crea una clave AES-256 unica
    \item \textbf{Derivacion HKDF}: La clave se procesa mediante HKDF-SHA256
    \item \textbf{Hash Argon2id}: Se aplica Argon2id para fortalecer la clave
    \item \textbf{Cifrado AES-256-GCM}: Los datos se cifran con autenticacion
    \item \textbf{Nonce}: Se genera un nonce de 96 bits unico por operacion
    \item \textbf{Almacenamiento}: Datos cifrados persistidos de forma segura
\end{enumerate}

\section{Topologia de Red}
\label{sec:topologia-red}

\begin{table}[H]
\centering
\begin{tabular}{p{4cm}p{4cm}p{5cm}}
\toprule
\textbf{Componente} & \textbf{Conexion} & \textbf{Protocolo} \\
\midrule
Nodos Cliente (SDC) & Internet & TLS 1.3 \\
\midrule
Servidor Principal & Internet & WSS (Puerto 8443) \\
\midrule
Servidor Replica & Internet & WSS (Puerto 8443) \\
\midrule
Base de Datos & Servidores & AES-256 \\
\bottomrule
\end{tabular}
\caption{Topologia de red del ecosistema SDC}
\label{tab:topologia-red}
\end{table}

\section{Tablas de Referencia Rapida}
\label{sec:tablas-ref}

\subsection{Puertos y Protocolos}

\begin{table}[H]
\centering
\begin{tabular}{p{2cm}p{3cm}p{4cm}p{4cm}}
\toprule
\textbf{Puerto} & \textbf{Protocolo} & \textbf{Uso} & \textbf{Seguridad} \\
\midrule
8443 & WSS & Comunicacion cliente-servidor & TLS 1.3 + Client Auth \\
\midrule
8080 & WS & Desarrollo local & Sin cifrado \\
\midrule
443 & HTTPS & API REST & TLS 1.3 \\
\midrule
3000 & HTTP & Desarrollo Angular & Solo localhost \\
\bottomrule
\end{tabular}
\caption{Puertos y protocolos utilizados por SDC}
\label{tab:puertos}
\end{table}

\subsection{Algoritmos Criptograficos}

\begin{table}[H]
\centering
\begin{tabular}{p{4cm}p{4cm}p{5cm}}
\toprule
\textbf{Proposito} & \textbf{Algoritmo} & \textbf{Parametros} \\
\midrule
Cifrado Simetrico & AES-256-GCM & Clave: 256 bits, Nonce: 96 bits \\
\midrule
Hashing de Contrasenas & Argon2id & Memoria: 64MB, Iteraciones: 3 \\
\midrule
Derivacion de Claves & HKDF-SHA256 & Salt: 128 bits \\
\midrule
Cifrado Asimetrico & RSA-4096 & OAEP padding, SHA-256 \\
\midrule
Firma Digital & ECDSA P-256 & Curva secp256r1 \\
\bottomrule
\end{tabular}
\caption{Algoritmos criptograficos implementados}
\label{tab:crypto}
\end{table}

\subsection{Codigos de Error}

\begin{table}[H]
\centering
\begin{tabular}{p{3cm}p{4cm}p{6cm}}
\toprule
\textbf{Codigo} & \textbf{Descripcion} & \textbf{Accion Recomendada} \\
\midrule
SDC-E001 & Fallo de autenticacion & Verificar credenciales y Client ID \\
\midrule
SDC-E002 & Conexion rechazada & Verificar red y estado del servidor \\
\midrule
SDC-E003 & Certificado invalido & Actualizar certificados CA del sistema \\
\midrule
SDC-E004 & Timeout agotado & Aumentar timeout o verificar latencia \\
\midrule
SDC-E005 & Violacion de seguridad & Contactar administrador inmediatamente \\
\midrule
SDC-E006 & Espacio insuficiente & Liberar espacio en disco \\
\midrule
SDC-E007 & Formato no soportado & Verificar compatibilidad del archivo \\
\midrule
SDC-E008 & Sesion expirada & Reiniciar sesion de autenticacion \\
\bottomrule
\end{tabular}
\caption{Codigos de error estandar de SDC}
\label{tab:errores}
\end{table}

\section{Glosario de Terminos Tecnicos Extendido}
\label{sec:glosario-tecnico}

\begin{description}[leftmargin=*, style=nextline]
    \item[\textbf{AES-256-GCM}] Advanced Encryption Standard con clave de 256 bits en modo Galois/Counter Mode. Proporciona cifrado autenticado.
    
    \item[\textbf{Argon2id}] Variante del algoritmo Argon2 disenada para resistir ataques de canal lateral y GPU.
    
    \item[\textbf{Client ID}] Identificador unico criptografico asignado a cada nodo SDC.
    
    \item[\textbf{Defensa en Profundidad}] Estrategia de seguridad que implementa multiples capas de proteccion.
    
    \item[\textbf{HKDF}] HMAC-based Key Derivation Function. Algoritmo para derivar multiples claves criptograficas.
    
    \item[\textbf{IPC}] Inter-Process Communication. Mecanismo de comunicacion entre Angular y Rust.
    
    \item[\textbf{Ruta Proxy}] Punto de acceso configurado que actua como intermediario seguro.
    
    \item[\textbf{Sandbox}] Entorno de ejecucion aislado donde las micro-aplicaciones operan.
    
    \item[\textbf{SID}] Secure ID. Identificador unico trazable asignado a cada operacion.
    
    \item[\textbf{WSS}] WebSocket Secure. Protocolo de comunicacion bidireccional sobre TLS.
\end{description}
